\section{本章小结}
\label{sec:vr-summary}

本章固定积分$I=\mathbb{E}_P[f(X)]$，研究怎样在可比的预算下减小随机计算的波动。比较的口径是第~\ref{sec:vr-comparison}节的单元方差乘以单元成本：一种改动只有在$v\,c$变小、而且计入预备成本后仍然划算时才算改进；有偏的估计比较均方误差，并单独检查偏差是否随样本量消失。

三类方法作用在计算流水线的不同位置。改变分布调整重要贡献被抽到的频率：重要性采样按建议分布抽样、乘权重折回，理想建议正比于贡献$\lvert f\rvert p$；它在覆盖条件下无偏，但方差有限要另查加权二阶矩，尾部过轻的建议会使方差为无穷；未知归一化常数时的自归一化估计只是一致的，有限样本有偏，权重的有效样本量看不到漏掉的区域。

改变计算公式消去可以预测或可以精确平均的波动：控制变量减去回归可解释的部分，条件期望把组内波动整个算掉。改变抽样机制控制点的去向与点之间的协方差：分层去掉层间波动，拉丁超立方去掉各坐标的主效应，随机化准蒙特卡洛按Fourier频率过滤误差，对偶变量利用负协方差，共同随机数利用正协方差。

\begin{BookInlineTable}
 \centering
 \caption{本章方法在贯穿例子上的效果。单样本方差均折合到一次函数求值；“倍数”是独立抽样方差与之的比值，未计入额外成本。}
 \label{tab:vr-summary}
 \begin{tabular}{llll}
  \toprule
  方法 & 起作用的条件 & 例子 & 倍数\\
  \midrule
  重要性采样 & $q$接近$\lvert f\rvert p$ & 尾概率，$\lambda=1/4$ & $10.3$\\
  试运行拟合 & 覆盖各区，留防御成分 & 两区域尾部 & $37.7$\\
  控制变量 & $\rho^2$接近$1$ & $\mathrm{e}^{U}$，$h=U$ & $61$\\
  条件期望 & 条件平均可算 & 四分之一圆 & $3.4$\\
  比例分层 & 层均值差别大 & $\mathrm{e}^{U}$，$4$层 & $14.6$\\
  随机移位 & $F$在环面上光滑 & $\mathrm{e}^{U}$，$4$点 & $3.9$\\
  对偶变量 & 奇部方差较大 & $\mathrm{e}^{U}$ & $31$\\
  共同随机数 & 两输出正相关 & 封顶等待时间 & $67.6$\\
  \bottomrule
 \end{tabular}
\end{BookInlineTable}

表~\ref{tab:vr-summary}中的倍数只在各自的例子上成立，并且每一行都有反例：$\lambda\geq2$的指数建议、没有防御成分的拟合、交互型函数上的拉丁超立方、对称函数上的对偶对、用反向输入的共同随机数，都会使方差不降反升甚至为无穷。生成映射$T$也会改变结论，同一个光滑的$f$经过不同的$T$，随机化准蒙特卡洛的误差速度可以从$n^{-2}$降到$n^{-2/3}$。

样本不再独立时，误差估计要相应改变：整组随机化的方法以完整的一组为独立单位，用组间方差除以重复次数估计误差，不能用组内点的样本方差代替。三类操作可以组合，例如在拉丁超立方的点上再用控制变量，或在分层的每层内用重要性采样；但期望正确、方差有限与节省总计算量是三个不同的要求，组合后要重新逐项检查，各方法单独使用时的改进不能直接相乘。实际预算下怎样分配试运行与正式计算、怎样在方法之间取舍，由第~\ref{chap:sampling-budget-method-selection}章统一讨论；下一章把同样的思路用于含参期望的梯度。
